% teorema-degrau2-draft.tex -- PROVA do TEOREMA DEGRAU 2 (readout treinavel).
% [TEOREMA -- promovido no 18o parecer (2 rodadas), 2026-07-27]. Status: (H6') sobre v_0
% como hipotese, verificavel a posteriori; (P1)-(P10) abertos declarados; beta_0 herda
% e^{X_max C_res H} para todas as clausulas.  NAO e o paper.  Fonte da verdade numerica:
%   06-notas/2026-07-27-degrau2-scoping.md
%   05-experimentos/t2-n-neuronios/degrau2_readout.{py,jsonl}   (v_0(0)=1)
%   05-experimentos/t2-n-neuronios/degrau2_v0.{py,jsonl}        (v_0(0) in {0.5,2})
%   06-notas/t2-prova/degrau2_lema_verificacao.py               (SymPy + numerico)
% Notacao do paper-v4.tex.  Compilar: pdflatex teorema-degrau2-draft.tex (2x).
\documentclass[11pt]{article}
\usepackage{amsmath,amssymb,amsthm}
\usepackage[text={6in,8.4in},centering]{geometry}
\usepackage{booktabs}
\usepackage{microtype}
\usepackage{url}
\usepackage[colorlinks=true,linkcolor=blue,citecolor=blue,urlcolor=blue]{hyperref}
\usepackage[nameinlink]{cleveref}

\newtheorem{theorem}{Theorem}
\newtheorem{lemma}[theorem]{Lemma}
\newtheorem{proposition}[theorem]{Proposition}
\newtheorem{corollary}[theorem]{Corollary}
\theoremstyle{remark}
\newtheorem{remark}[theorem]{Remark}
\crefname{theorem}{Theorem}{Theorems}
\crefname{lemma}{Lemma}{Lemmas}
\crefname{remark}{Remark}{Remarks}
\crefname{proposition}{Proposition}{Propositions}
\crefname{corollary}{Corollary}{Corollaries}
\crefname{section}{Section}{Sections}
\crefname{equation}{}{}
\Crefname{equation}{}{}

\newcommand{\R}{\mathbb{R}}
\newcommand{\thit}{t_{\mathrm{hit}}}
\newcommand{\texit}{t_{\mathrm{exit}}}
\newcommand{\trel}{t_{\mathrm{release}}}
\newcommand{\Ar}{\mathcal{A}}
\newcommand{\Tep}{T_{\mathrm{ep}}}
\newcommand{\hyb}{\mathrm{hyb}}
\newcommand{\res}{\mathrm{res}}
\newcommand{\Cres}{C_{\mathrm{res}}}
\newcommand{\Xmax}{X_{\max}}
\newcommand{\LipR}{\mathrm{Lip}_R}
\newcommand{\hyplabel}[1]{\textup{(#1)}}

\title{The Degrau~2 theorem: the entry--exit law with a \emph{trainable} readout\\
\large [THEOREM --- promoted at the 18th referee report (two rounds), 2026-07-27]}
\author{j-space internal note}
\date{2026-07-27}

\begin{document}
\maketitle

\begin{abstract}
Theorem~6 of \texttt{paper-v4} (the \emph{forcing} form of the entry--exit law)
assumes the live subsystem $B$ is \emph{affine of order $m$}: frozen activation
pattern \emph{and fixed readouts}. Open item (O6) of that paper asks what
survives when the readouts $v_k$ are trained too. This note answers: everything
survives, at the price of \textbf{one} new elementary lemma
(\cref{lem:curv}) and \textbf{one} new hypothesis \hyplabel{H6$'$} --- a guard
on the \emph{dead} unit's readout alone, no live readout being constrained
(\cref{rem:H6only0}) --- and with the
first-order coefficient acquiring an explicit and \emph{measurable} factor,
$-1/\bigl(c\,v_0(\thit)\,R^{\hyb}(\texit)\bigr)$. The new lemma says that the
negative curvature of the loss over the live block is bounded by
$\Lambda=\max_k\|c_k\|=\max_k\|\nabla_{(w_k,b_k)}L\|/|v_k|\le\Xmax\Cres$: a
\emph{maximum over blocks}, hence independent of $m$. This corrects (O6), which
predicts a Gr\"onwall factor $e^{L_B\Tep}$ with $L_B=\lambda_{\max}(G_B)\sim m$;
the correct exponent admits a bound free of the width. We prove the lemma
(verified symbolically in SymPy and numerically on $2\cdot10^4$ random points),
rewrite L2$'$ as L2$''$, and settle the coefficient by a \emph{discriminating}
numerical experiment with $v_0(0)\in\{0.5,2.0\}$ that the scoping note left
undone --- the previous campaign had $v_0(0)=1$ and was a null of instrument.
The experiment is decisive and it decides \emph{for} the corrected law. We are
explicit about what the exponential factor costs and, in particular, correct the
scoping note on where it enters: clause~(ivB) pays a \emph{horizon} factor, not
an episode factor, and the threshold $\beta_0$ inherits that horizon factor for
\emph{every} clause even though the statements of (i)--(iii) do not.
\end{abstract}

\tableofcontents

%======================================================================
\section{Setup, and the one new hypothesis}\label{sec:setup}

\subsection{The network and the two flows}

Fix data $(x_i,y_i)_{i=1}^n\subset\R\times\R$, write $\Xmax=\max_i\|(x_i,1)\|$,
and let $g_\beta(z)=\beta^{-1}\log(1+e^{\beta z})$ be the softplus, $\sigma$ the
logistic. The network is
\begin{equation}\label{eq:net}
f_\theta(x)=v_0\,g_\beta(w_0x+b_0)+\sum_{k=1}^{m}v_k\,g_\beta(w_kx+b_k),
\end{equation}
and --- this is the only change from \cite[Thm.~6]{paper} --- \textbf{all}
$3(m+1)$ parameters are trained:
$\theta=(v_k,w_k,b_k)_{k=0}^m\in\R^{3(m+1)}$. Unit $0$ is the distinguished
unit $A$ with margin $z=w_0x_j+b_0$ on the face $j$; the remaining $m$ units
form the live block $B$, $\theta_B=(v_k,w_k,b_k)_{k=1}^m\in\R^{3m}$. The loss
is $L(\theta)=\frac1{2n}\sum_i\res_i^2$, $\res_i=f_\theta(x_i)-y_i$, the flow is
$\dot\theta=-\nabla L(\theta)$, and
\[
R(t)=y_j-f_{-A}(x_j;\theta_B(t)),\qquad c=\frac{x_j^2+1}{n},
\]
is the forcing generated by $B$ at the face. The \emph{hybrid} reference
$\theta^{\hyb}$ is as in \cite[Def.~2]{paper}: ReLU dynamics, with the whole
triple $(v_0,w_0,b_0)$ frozen at the value where $z=0$ from $\thit$ onwards;
$\theta^\infty$ is the hybrid limit in which $A$ revives at $\texit$. Put
$\Delta=\theta_B^\beta-\theta_B^{\hyb}$, $\rho=\|\Delta\|$,
$\Tep=\texit-\thit$.

\subsection{What breaks, precisely}\label{sec:breaks}

With the readouts frozen, $f_{-A}$ is \emph{linear} in $\theta_B$ on a window of
constant activation pattern: $f_{-A}(x_i;\theta_B)=\varphi_i\cdot\theta_B$, the
ReLU $B$-field is $-G_B(\theta_B-\theta_B^\ast)$ with $G_B=\frac1n\sum_i
\varphi_i\varphi_i^\top\succeq0$, and step~1 of \cite[L2$'$]{paper} is the exact
identity $\dot\Delta=-G_B\Delta+G_A+G_S$ --- the two flows share the \emph{same}
matrix, so the difference of drifts has no remainder. Step~3 then reads
$\frac{d}{dt}\frac12\rho^2=-\Delta^\top G_B\Delta+\Delta^\top(G_A+G_S)
\le\rho(\|G_A\|+\|G_S\|)$ and \emph{no exponential is paid}. That single sign is
what makes $M_2$ polynomial in $m$.

With trainable readouts and pattern frozen,
\begin{equation}\label{eq:bilinear}
f_{-A}(x_i;\theta_B)=\sum_{k=1}^m v_k\,\mathbf1_{k,i}\,(w_kx_i+b_k)
\end{equation}
is \emph{bilinear}, not linear. There is no fixed $G_B$; the affine class of
\cite[\S3.1]{paper} is gone, and with it the literal form of step~1.

But what step~1 uses is not affineness --- it is that \emph{both flows run the
same vector field}. Fix the (constant, by \hyplabel{H4$'$}) indicators
$\mathbf1_{k,i}$ and let
\begin{equation}\label{eq:Ufrozen}
U(\theta_B):=\frac{1}{2n}\sum_{i=1}^n
\Bigl(\textstyle\sum_{k}v_k\mathbf1_{k,i}(w_kx_i+b_k)-y_i\Bigr)^2
\end{equation}
be the frozen-pattern surrogate potential: a \emph{polynomial}, hence $C^\infty$
on all of $\R^{3m}$, with no kinks to worry about. The hybrid $B$-flow is
exactly $-\nabla U$; the $\beta$-flow is $-\nabla U$ plus two sources ($A$'s
output, and the softplus defects of the live units). By the mean value theorem
applied to the gradient,
\begin{equation}\label{eq:Deltaode}
\boxed{\ \dot\Delta=-\bar H(t)\,\Delta+G_A(t)+G_S(t),\qquad
\bar H(t)=\int_0^1\nabla^2U\bigl(\theta_B^{\hyb}(t)+s\Delta(t)\bigr)\,ds\ }
\end{equation}
with $\bar H(t)$ \emph{symmetric}. This is still an identity, not an estimate.
Exactly one thing has been lost: $\bar H\succeq0$. \Cref{lem:curv} recovers as
much of it as is needed.

\begin{remark}[why \eqref{eq:Deltaode} has no pattern subtlety]\label{rem:noswitch}
The interpolating segment $\theta_B^{\hyb}+s\Delta$ can leave the region where
the true ReLU pattern is the frozen one. This does not matter: $U$ in
\eqref{eq:Ufrozen} is defined with the indicators as \emph{constants}, so it is
a globally smooth polynomial and the mean value theorem applies unconditionally.
The discrepancy between $U$ and the true ReLU potential off the pattern region
never enters, because \emph{both} flows are compared to $U$, the hybrid exactly
(by \hyplabel{H4$'$}) and the $\beta$-flow up to $G_S$. This is a genuine
simplification relative to the situation of (O2) in \cite{paper}, where the
\emph{flows themselves} switch at different instants.
\end{remark}

\subsection{Hypotheses}

We assume \hyplabel{H0$'$}--\hyplabel{H5$'$} of \cite[\S3.1]{paper} verbatim,
with two readings made explicit: (i) in \hyplabel{H1$'$}/\hyplabel{H2$'$} the
arrival and exit speeds are now $v=c\,|v_0(\thit)|\,|R^{\hyb}(\thit)|$ and
$v'=c\,v_0(\thit)\,R^{\hyb}(\texit)$ (\cref{sec:coef}); (ii) \hyplabel{H5$'$}
($|\res_i|\le\Cres$) and \hyplabel{H4$'$} (uniform live margins $\ge\delta_B$)
are required \emph{on the tube}, i.e.\ on the segments
$\{\theta_B^{\hyb}+s\Delta:s\in[0,1]\}$, exactly as \cite{paper} already states
them (``along both trajectories in the tube''). We add one hypothesis.

\begin{itemize}
\item[\hyplabel{H6$'$}] \textbf{Guard on the dead unit's readout.} There is
$\varphi_v>0$ with
\[
\min_{t\in[0,H]}\ \bigl|v_0(t)\bigr|\ \ge\ \varphi_v>0,
\]
and $v_0$ does not change sign on $[0,H]$, along both trajectories and on the
tube. \textbf{Nothing is assumed about the live readouts $v_1,\dots,v_m$.}
\end{itemize}

\begin{remark}[why the guard is needed on $v_0$ only]\label{rem:H6only0}
The natural first draft of \hyplabel{H6$'$} asks $\min_{0\le k\le
m}\min_t|v_k|\ge\varphi_v$. That is stronger than anything the proof consumes,
and we weaken it deliberately. Audit of every place a readout appears:
\begin{itemize}
\item $\Lambda$ of \cref{lem:curv} is $\max_k\|c_k\|$ with
$c_k=\frac1n\sum_i\res_i\mathbf1_{k,i}(x_i,1)$: \emph{it does not mention $v_k$
at all}, and the bound used, $\Lambda\le\Xmax\Cres$, is free of readouts. The
identity $\Lambda=\max_k\|\nabla_{(w_k,b_k)}U\|/|v_k|$ is a \emph{reading}, not
a step (\cref{rem:readings}).
\item $C_f$, $\|G_{BA}\|$, $C_5$ are \emph{upper} bounds on readout norms
(\cref{sec:unif}, U3, U-new). Upper bounds do not need a floor.
\item the arrival speed $v=c|v_0||R^{\hyb}(\thit)|$, the exit speed
$v'=cv_0R^{\hyb}(\texit)$, the constant $C_4$ and the coefficient
\eqref{eq:release} involve $v_0$ and \emph{only} $v_0$.
\end{itemize}
So a live readout may vanish, or change sign, without any step of
L1$''$--L6$''$ failing: the live block simply loses a direction of its own
Gauss--Newton part, which is a PSD term we only ever discard. See
\cref{sec:open}\,(P10).
\end{remark}

\begin{remark}[\hyplabel{H6$'$} is what Algorithm~1 already enforces]
\hyplabel{H6$'$} is not a new demand on the experimenter: the acceptance filter
of \texttt{degrau2\_readout.py} (the trainable-readout descendant of
Algorithm~1 of \cite{paper}) rejects, instance by instance, any run in which
\emph{some} readout crosses zero inside the window, on \emph{either} the hybrid
or the $\beta$-run; that filter is \emph{stricter} than \hyplabel{H6$'$} as now
stated, so the accepted instances satisfy it a fortiori, and $\varphi_v$ is read
off as $\min_t|v_0|$ over the accepted trajectory. So \hyplabel{H6$'$} is checkable a posteriori by the same
single integration that checks \hyplabel{H1$'$}--\hyplabel{H2$'$}, which is the
standard this paper holds all its hypotheses to.

Under ReLU gradient flow each unit conserves exactly
$B_k=v_k^2-\|(w_k,b_k)\|^2$ (Du--Hu--Lee balancedness, \cite{duhulee}), so
$B_k(0)>0$ gives $|v_k(t)|\ge\sqrt{B_k(0)}$ for all $t$ and makes the guard a
condition on the \emph{initialization}.

\textbf{Corollary of the reduction to $k=0$.} Because \hyplabel{H6$'$} now
concerns one unit, the balancedness route needs
\[
B_0(0)=v_0(0)^2-\|(w_0,b_0)(0)\|^2>0
\]
for the \emph{dead unit alone}, and then $\varphi_v=\sqrt{B_0(0)}$. That is a
condition on \emph{one} triple, imposable by construction --- e.g.\ by the
balanced rescaling $(v_0,w_0,b_0)\mapsto(sv_0,w_0/s,b_0/s)$, which does not
change $f$ --- whereas the $m+1$-unit version competed with \hyplabel{H4$'$} on
every live unit and failed (\cref{sec:open}\,(P3)). The campaigns here still use
the initialization of \cite{paper}, which has $B_0(0)<0$, so we \emph{measure}
$\varphi_v$ rather than construct it; but the route from hypothesis to
construction is now a one-unit problem, and that is the substantive gain.
\end{remark}

%======================================================================
\section{The new lemma: block-diagonal negative curvature}\label{sec:lema}

This is the whole mathematical content of Degrau~2. It is elementary linear
algebra and we write it out in full.

\begin{lemma}[block-diagonal negative curvature]\label{lem:curv}
Let $U$ be as in \eqref{eq:Ufrozen} (one hidden layer, frozen activation
pattern, quadratic loss, trainable readouts). Write, for $k=1,\dots,m$,
\begin{equation}\label{eq:ck}
c_k(\theta_B):=\frac1n\sum_{i=1}^n \res_i\,\mathbf1_{k,i}\,(x_i,1)\ \in\R^2,
\qquad \res_i=\sum_{l}v_l\mathbf1_{l,i}(w_lx_i+b_l)-y_i .
\end{equation}
Then, at every $\theta_B\in\R^{3m}$,
\begin{equation}\label{eq:curv}
\lambda_{\min}\bigl(\nabla^2U(\theta_B)\bigr)\ \ge\ -\Lambda(\theta_B),
\qquad
\Lambda(\theta_B):=\max_{1\le k\le m}\|c_k(\theta_B)\| ,
\end{equation}
and, whenever $v_k\neq0$ for all $k$,
\begin{equation}\label{eq:twoReadings}
\Lambda=\max_{1\le k\le m}\frac{\|\nabla_{(w_k,b_k)}U\|}{|v_k|}
\qquad\text{and}\qquad
\Lambda\ \le\ \Xmax\,\Cres .
\end{equation}
\end{lemma}

\begin{proof}
Order the coordinates \emph{by unit}: $\theta_B=(w_1,b_1,v_1,\dots,w_m,b_m,v_m)$.

\smallskip
\emph{Step 1 (Gauss--Newton splitting).} With
$f_i(\theta_B)=\sum_kv_k\mathbf1_{k,i}(w_kx_i+b_k)$ and
$\res_i=f_i-y_i$, differentiating $U=\frac1{2n}\sum_i\res_i^2$ twice gives the
exact identity
\begin{equation}\label{eq:GNsplit}
\nabla^2U
=\underbrace{\frac1n\sum_{i=1}^n\nabla f_i\,\nabla f_i^{\top}}_{=\;\frac1nJ^\top J\;\succeq\;0}
\;+\;\underbrace{\frac1n\sum_{i=1}^n \res_i\,\nabla^2f_i}_{=:\,S},
\end{equation}
$J\in\R^{n\times3m}$ being the Jacobian $J_{i\cdot}=\nabla f_i^\top$. The first
term is a Gram matrix, hence PSD, for every $\theta_B$ and every $m$. It is the
generalization of $G_B$: with the $v_k$ frozen it \emph{is} $G_B$.

\smallskip
\emph{Step 2 ($S$ is block diagonal, $3\times3$ per unit).} Each $f_i$ is a sum
of terms $v_k\mathbf1_{k,i}(w_kx_i+b_k)$, each involving the parameters of a
\emph{single} unit and being \emph{bilinear} in $(v_k)$ versus $(w_k,b_k)$.
Hence all second derivatives across distinct units vanish,
$\partial^2f_i/\partial\theta_k\partial\theta_l=0$ for $k\ne l$, and within
unit $k$ the only nonvanishing ones are the mixed pair
\[
\frac{\partial^2f_i}{\partial v_k\,\partial w_k}=\mathbf1_{k,i}\,x_i,
\qquad
\frac{\partial^2f_i}{\partial v_k\,\partial b_k}=\mathbf1_{k,i},
\qquad
\frac{\partial^2f_i}{\partial v_k^2}
=\frac{\partial^2f_i}{\partial w_k^2}
=\frac{\partial^2f_i}{\partial w_k\partial b_k}
=\frac{\partial^2f_i}{\partial b_k^2}=0 .
\]
Therefore $S=\mathrm{diag}(S_1,\dots,S_m)$ with, in the order $(w_k,b_k,v_k)$,
\begin{equation}\label{eq:Sk}
S_k=\begin{pmatrix}0&0&c_{k,1}\\ 0&0&c_{k,2}\\ c_{k,1}&c_{k,2}&0\end{pmatrix}
=\begin{pmatrix}0&c_k\\ c_k^\top&0\end{pmatrix},
\qquad c_k=(c_{k,1},c_{k,2})=\frac1n\sum_i\res_i\mathbf1_{k,i}(x_i,1),
\end{equation}
which is \eqref{eq:ck}. \emph{This is the load-bearing structural fact}: the
residual curvature never couples two different neurons, because a one-hidden-layer
network is a \emph{sum} over units and each summand is a function of one unit's
parameters alone.

\smallskip
\emph{Step 3 (spectrum of a $3\times3$ block).} Fix $k$ and write $c=c_k$. If
$c=0$ then $S_k=0$. Otherwise, the characteristic polynomial is
\[
\det(S_k-\mu I)
=-\mu^3+\mu\,(c_1^2+c_2^2)
=-\mu\bigl(\mu^2-\|c\|^2\bigr),
\]
so $\mathrm{spec}(S_k)=\{0,+\|c\|,-\|c\|\}$. (Directly: $S_k$ has rank $2$, its
kernel is spanned by $(-c_2,c_1,0)$; on the invariant plane spanned by
$e_c:=(c_1,c_2,0)/\|c\|$ and $e_3=(0,0,1)$ it acts as
$\|c\|\left(\begin{smallmatrix}0&1\\1&0\end{smallmatrix}\right)$, with eigenvalues
$\pm\|c\|$ and eigenvectors $(e_c\pm e_3)/\sqrt2$.) Hence
$\lambda_{\min}(S_k)=-\|c_k\|$ and, $S$ being block diagonal,
\begin{equation}\label{eq:lamS}
\lambda_{\min}(S)=\min_{1\le k\le m}\lambda_{\min}(S_k)
=-\max_{1\le k\le m}\|c_k\|=-\Lambda .
\end{equation}
The $\max$ in \eqref{eq:lamS} is the entire point. Had the residual curvature
coupled the units, the estimate would have run through a trace or a sum and
picked up a factor $m$ (or $\sqrt m$); because $S$ is block diagonal, the
spectrum of $S$ is the \emph{union} of the block spectra and the extreme
eigenvalue is a \emph{maximum over $m$ numbers}, each of which is bounded
independently of $m$.

\emph{A precision the referee insisted on.} What is $m$-free is the
\textbf{bound} $\Lambda\le\Xmax\Cres$, which is what every estimate below uses.
The \emph{realized} $\Lambda$ along a trajectory is an empirical quantity and
could in principle drift with $m$ through the residuals. A pilot (accepted
hybrid instances, coarse $dt=2\cdot10^{-4}$, script \path{parecer18_m2.py})
finds median $\Lambda=5.03,\,5.70,\,5.18$ at $m=4,8,16$ against
$\lambda_{\max}(G_B)=10.2,\,18.7,\,37.7$ --- the latter plainly linear in $m$,
the former flat. $N$ is $4$, $4$ and $1$ at the three levels after filtering, so
this is \emph{indicative and nothing more}; the proper sweep is
\cref{sec:open}\,(P7).

\smallskip
\emph{Step 4 (Weyl).} By Weyl's inequality applied to \eqref{eq:GNsplit},
\[
\lambda_{\min}(\nabla^2U)\ \ge\ \lambda_{\min}\bigl(\tfrac1nJ^\top J\bigr)+\lambda_{\min}(S)
\ \ge\ 0+(-\Lambda)=-\Lambda,
\]
which is \eqref{eq:curv}.

\smallskip
\emph{Step 5 (the two readings).} From \eqref{eq:Ufrozen},
$\nabla_{(w_k,b_k)}U=\frac1n\sum_i\res_i\,v_k\mathbf1_{k,i}(x_i,1)=v_k\,c_k$,
whence $\|c_k\|=\|\nabla_{(w_k,b_k)}U\|/|v_k|$ when $v_k\ne0$: the first half of
\eqref{eq:twoReadings}. For the second, $\mathbf1_{k,i}\in\{0,1\}$ and the
triangle inequality give
$\|c_k\|\le\frac1n\sum_i|\res_i|\,\|(x_i,1)\|\le\Cres\Xmax$, uniformly in $k$
and in $m$.
\end{proof}

\begin{remark}[the two readings, and which one is used where]\label{rem:readings}
\emph{Crude:} $\Lambda\le\Xmax\Cres$ involves only the data and the residual
bound of \hyplabel{H5$'$}. It carries \textbf{no $m$}, no $\varphi_v$, and no
readout norm --- observe that $\Lambda$ as defined in \eqref{eq:ck} does not
mention $v_k$ at all. This is the bound used in every estimate below.
\emph{Fine:} $\Lambda=\max_k\|\nabla_{(w_k,b_k)}U\|/|v_k|\le
\|\nabla_{\theta_B}U\|/\varphi_v$ decays as the flow converges, and is the
reading in which the guard \hyplabel{H6$'$} appears. It is sharper late in
training and it is what one would use to show that $\int_0^H\Lambda\,dt\ll
\Lambda_\infty H$; we do not need it for the theorem, and we do not claim it.
\end{remark}

\begin{remark}[this corrects (O6)]\label{rem:O6}
Item (O6) of \cite{paper} predicts that a trainable readout brings back the
Gr\"onwall factor $e^{L_B\Tep}$ with $L_B=\lambda_{\max}(G_B)$. The $\Tep$ is
right; the $L_B$ is not. \emph{The exponent is not the largest eigenvalue of the
PSD part but the most negative eigenvalue of the residual part}, and those are
different objects: $\lambda_{\max}(G_B)$ grows like $m$ (it is a sum of $n$ rank-one
terms whose norms add over units), whereas $\Lambda$ is a maximum over blocks and
does not. As stated, (O6) is more pessimistic than the mathematics requires,
and \cref{lem:curv} is the correction.

\emph{Notation, to keep the two apart.} We reserve $L_B$ for
$\lambda_{\max}(G_B)$, the paper's object, and write
$\Lambda_{\rm meas}:=\max\bigl(0,-\lambda_{\min}(\nabla^2U)\bigr)$ for the
quantity the experiments measure and that \cref{lem:curv} bounds by $\Lambda$.
Using one symbol for both --- as an earlier draft of this note did --- makes the
present remark unreadable, since its entire content is that they differ.
\end{remark}

\begin{remark}[executable verification]\label{rem:verif}
\cref{lem:curv} is verified by
\path{06-notas/t2-prova/degrau2_lema_verificacao.py} (run
\texttt{python degrau2\_lema\_verificacao.py}), in three independent ways.
\textbf{(A)} SymPy computes $\mathrm{spec}(S_k)$ for the symbolic block
\eqref{eq:Sk} and confirms $\{0,+\sqrt{c_1^2+c_2^2},-\sqrt{c_1^2+c_2^2}\}$
(step~3). \textbf{(B)} SymPy builds the \emph{full} symbolic Hessian of
\eqref{eq:Ufrozen} for $m=2$, $n=2$ with symbolic
$w,b,v,x,y$ and symbolic indicators, and verifies the identity
$\nabla^2U-\frac1nJ^\top J=\mathrm{diag}(S_1,S_2)$ \emph{exactly} --- in
particular that the inter-neuron residual coupling is identically zero
(steps~1--2). \textbf{(C)} $2\cdot10^4$ random points with $m\in[1,6]$,
$n\in[1,8]$: the relative slack $(\lambda_{\min}+\Lambda)/\Lambda$ has minimum
$-9.6\cdot10^{-14}$ (machine zero, i.e.\ the bound never fails) and is below
$10^{-10}$ --- i.e.\ \emph{saturated} --- in $22.9\%$ of points (step~4).
\end{remark}

\begin{remark}[saturation is observed, not proved]\label{rem:sat}
On the actual trajectories of \cref{sec:num}, \eqref{eq:curv} holds with
\emph{equality to machine precision}: over the $43$ instances of the
$v_0(0)=1$ campaign, $\max|\Lambda_{\rm meas}-\Lambda|=4.4\cdot10^{-15}$
(relative $9\cdot10^{-16}$), where
$\Lambda_{\rm meas}:=\max(0,-\lambda_{\min}(\nabla^2U))$ is measured by
exact eigenvalues of the $12\times12$ Hessian. The sharp form of the mechanism,
as the referee put it: Weyl closes with equality exactly when the
\emph{minimizing eigenvector of $S$ is also a null vector of the Gauss--Newton
part}. By step~3 of \cref{lem:curv} that eigenvector is explicit --- inside the
binding block $k^\ast=\arg\max_k\|c_k\|$ it is $(e_{c_{k^\ast}}-e_3)/\sqrt2$,
and it is zero on every other block. Since
$\mathrm{rank}(\tfrac1nJ^\top J)\le n<3m$ (here $4<12$), the Gauss--Newton
kernel has dimension $\ge3m-n=8$; the claim is that this \emph{particular}
vector lands in it. That is \emph{not} automatic from the rank deficit: the random-point test of
\cref{rem:verif}\,(C) has the same deficit and saturates only $23\%$ of the
time. \textbf{What is a theorem is the inequality \eqref{eq:curv}; the equality
is an empirical fact about this family.} [SPECULATION for the mechanism.]
\end{remark}

%======================================================================
\section{\texorpdfstring{L2$''$}{L2''}: the live subsystem with an averaged Hessian}\label{sec:L2}

\begin{lemma}[live subsystem and feedback, averaged-Hessian form]\label{lem:L2}
Assume \hyplabel{H4$'$}, \hyplabel{H5$'$}, \hyplabel{H6$'$} and, on $[\thit,s]$,
the layer bound $\int_{\thit}^{s}c\,v_0\log(1+w)\,dt\le2C_4\log\beta/\beta$
supplied by the bootstrap L4$''$. Put
\begin{equation}\label{eq:Lamsup}
\Lambda_{\rm ep}:=\sup_{t\in[\thit,\,s]}\ \sup_{\sigma\in[0,1]}
\Lambda\bigl(\theta_B^{\hyb}(t)+\sigma\Delta(t)\bigr)\ \le\ \Xmax\Cres ,
\end{equation}
the supremum being taken over the \emph{same} interval $[\thit,s]$ on which the
bootstrap hypothesis is assumed, not over a fixed $[\thit,\trel]$; the induction
of L4$''$ closes on $s$ and the two must match. Then
\begin{multline}\label{eq:Brate}
\sup_{[\thit,s]}\rho
\ \le\ e^{\Lambda_{\rm ep}(s-\thit)}
\Bigl[\underbrace{\rho(\thit)}_{\text{initial datum}}
+\underbrace{\frac{C_fC_4}{c\,|v_0|}\cdot\frac{2\log\beta}{\beta^2}}_{\text{episode}}\\
{}+\underbrace{H\bigl(C_3e^{-\beta\delta_1}+C_5\sqrt m\,\Cres e^{-\beta\delta_B}\bigr)}_{\text{beyond all orders}}\Bigr],
\qquad
\rho(\thit)\le O(\beta^{-2})\,e^{\Lambda_{\rm pre}\thit},
\end{multline}
so that on $[\thit,\trel]$, where $s-\thit=\Tep+O(\beta^{-1})$, the prefactor is
$e^{\Lambda_{\rm ep}\Tep}(1+O(\beta^{-1}))$ and the pre-hit factor
$e^{\Lambda_{\rm pre}\thit}$ sits \emph{inside} the bracket --- the two
\textbf{compose}; they are not two independent readings of one factor ---
and $\sup\rho=O(\log\beta/\beta^2)$ with
$M_2^{\,\prime\prime}=e^{\Lambda_{\rm ep}\Tep}\,C_fC_4/(c\,|v_0|)$, and
consequently $R^\beta=R^{\hyb}+\delta$ with
$\|\delta\|_\infty\le\LipR\sup\rho+m\beta^{-1}e^{-\beta\delta_B}$.
\emph{$\lambda_{\max}$ of the Gauss--Newton part never appears; the exponent is
$\Lambda\le\Xmax\Cres$, independent of $m$.}
\end{lemma}

\begin{proof}
\emph{Step 1 (exact perturbed system).} This is \eqref{eq:Deltaode}, established
in \cref{sec:breaks} and \cref{rem:noswitch}. The sources are unchanged from
\cite[L2$'$]{paper}:
\[
G_A(t)=-\frac1n\sum_i v_0\,g_\beta(z_{A,i}(t))\,\nabla_{\theta_B}f_i,
\qquad G_S=\text{live softplus defect},
\]
the only novelty being the explicit factor $v_0$ in $G_A$ (the dead unit's
contribution to every residual is $v_0g_\beta(z_{A,i})$, not $g_\beta(z_{A,i})$)
and the fact that $\nabla_{\theta_B}f_i$ now has a third component per unit,
$\mathbf1_{k,i}(w_kx_i+b_k)$. Both are absorbed by redefining
$C_f:=\sup_t\|\nabla_{\theta_B}f_j\|/n$, still $O(\sqrt m)$ (\cref{sec:unif}).

\emph{Steps 2a, 2b (the sources) are untouched.} Step~2a bounds $G_A$ using only
\hyplabel{H4$'$} (the other margins of $A$ are $\le-\delta_1$) and yields
$\|G_A\|\le C_f|v_0|\beta^{-1}\log(1+w)+C_3e^{-\beta\delta_1}$. Step~2b bounds
$G_S$ using only \hyplabel{H4$'$} ($|z_{k,i}|\ge\delta_B$) and \hyplabel{H5$'$},
yielding $\|G_S\|\le C_5\sqrt m\,\Cres e^{-\beta\delta_B}$ for
$\beta\ge\beta_2$. Neither step says anything about the \emph{class} of the
live subsystem --- they are statements about $A$'s output and about the
softplus-versus-ReLU discrepancy of the live units --- so both survive verbatim.
This is worth stating because it is where one might have expected trouble and
there is none.

\emph{Step 3 (energy estimate, now with a Gr\"onwall).} Let $\rho=\|\Delta\|$.
From \eqref{eq:Deltaode},
\[
\frac{d}{dt}\tfrac12\rho^2
=-\Delta^\top\bar H(t)\Delta+\Delta^\top(G_A+G_S)
\ \le\ \Lambda_{\rm ep}\,\rho^2+\rho\,(\|G_A\|+\|G_S\|),
\]
where the first inequality is \cref{lem:curv} applied pointwise inside the
average defining $\bar H$: for each $s\in[0,1]$,
$-\Delta^\top\nabla^2U(\theta_B^{\hyb}+s\Delta)\Delta\le
\Lambda(\theta_B^{\hyb}+s\Delta)\rho^2\le\Lambda_{\rm ep}\rho^2$, and averaging
over $s$ preserves the bound. Wherever $\rho>0$ this gives
$\dot\rho\le\Lambda_{\rm ep}\rho+\|G_A\|+\|G_S\|$; $\rho$ is Lipschitz, so the
same holds for its upper Dini derivative everywhere, and Gr\"onwall's inequality
in integral form yields, for $t\in[\thit,s]$,
\begin{equation}\label{eq:gronwall}
\rho(t)\ \le\ e^{\Lambda_{\rm ep}(t-\thit)}
\Bigl[\rho(\thit)+\int_{\thit}^{t}e^{-\Lambda_{\rm ep}(u-\thit)}
\bigl(\|G_A(u)\|+\|G_S(u)\|\bigr)\,du\Bigr].
\end{equation}
Bounding $e^{-\Lambda_{\rm ep}(u-\thit)}\le1$ and substituting the step-2 bounds
together with the layer bound of the hypothesis gives \eqref{eq:Brate}, with the
prefactor evaluated at $t=\trel$: by clause~(ii),
$\trel-\thit=\Tep+O(\beta^{-1})$, so
$e^{\Lambda_{\rm ep}(\trel-\thit)}=e^{\Lambda_{\rm ep}\Tep}(1+O(\beta^{-1}))$.
\emph{This is the only place where $e^{\Lambda\Tep}$, and $\Tep$, occur.}

\emph{Step 4 (pre-hit gap $\rho(\thit)$; the factor here is \textbf{not}
$e^{\Lambda\Tep}$).} On $[0,\thit]$ both flows have $A$ alive with frozen
pattern, so the same argument applies to the \emph{full} vector
$\delta\theta=\theta^\beta-\theta^{\rm ReLU}\in\R^{3(m+1)}$ with the full frozen
potential; \cref{lem:curv} applies verbatim to the full network (the block
decomposition \eqref{eq:Sk} runs over $k=0,\dots,m$), giving
$\lambda_{\min}\ge-\Lambda_{\rm pre}$ with
$\Lambda_{\rm pre}=\sup_{[0,\thit]}\max_{0\le k\le m}\|c_k\|\le\Xmax\Cres$.
Step~4 of \cite[L2$'$]{paper} then goes through with the free PSD estimate
(B.24) of that paper replaced by \eqref{eq:gronwall}, and returns
\[
\rho(\thit)\ \le\ e^{\Lambda_{\rm pre}\thit}
\Bigl[\frac{C\|G_{BA}\|}{(\beta v)^2}+\frac{C_f}{\beta^2v}
+HC_5\sqrt m\,\Cres e^{-\beta\delta_B}\Bigr]
\ =\ O(\beta^{-2})\,e^{\Lambda_{\rm pre}\thit}.
\]
\textbf{The interval is $[0,\thit]$, of length $\thit$, and $\thit$ has nothing
to do with $\Tep$}; in the campaigns of \cref{sec:num} one typically has
$\thit\gg\Tep$. The rate is unaffected --- $O(\beta^{-2})$ times a constant is
still $O(\beta^{-2})$, and is dominated by the episode term, which carries an
extra $\log\beta$ --- but the \emph{constant} in the pre-hit group gains
$e^{\Lambda_{\rm pre}\thit}$, not $e^{\Lambda\Tep}$. The separation of one power
of $\beta$ used by (O5) of \cite{paper} does not use $\bar H\succeq0$, only the
sign of the defect, and is unaffected.

\emph{Step 5 (after release; \textbf{a horizon factor, not an episode factor}).}
On $[\trel,H]$ both the $\beta$-flow and $\theta^\infty$ have $A$ alive with the
same frozen pattern, so with $\delta\theta=\theta^\beta-\theta^\infty$ over the
full $\R^{3(m+1)}$ we again have the identity
$\dot{\delta\theta}=-\bar H_{\rm full}\,\delta\theta+E$ and again
$\lambda_{\min}(\nabla^2U_{\rm full})\ge-\Lambda_{\rm post}$ by \cref{lem:curv}
with $\Lambda_{\rm post}=\sup_{[\trel,H]}\max_{0\le k\le m}\|c_k\|\le\Xmax\Cres$
(sup over the tube). Repeating step~3 on $[\trel,H]$,
\begin{equation}\label{eq:ivB}
\sup_{[\trel,H]}\bigl\|\theta_B^\beta-\theta_B^{\infty}\bigr\|
\ \le\ \frac{M_3^{\,\prime\prime}}{\beta},
\qquad
M_3^{\,\prime\prime}=e^{\Lambda_{\rm post}(H-\trel)}
\Bigl[H\|G_{BA}\|C_A+2M_2^{\,\prime\prime}\tfrac{\log\beta}{\beta}+\tfrac{C_fH}{v'}\Bigr],
\end{equation}
with $v'=c\,v_0(\thit)R^{\hyb}(\texit)$ the exit speed and $C_A$ as in
\cite[(B.30)]{paper} with $c\mapsto c\,v_0$.
\end{proof}

\begin{remark}[where the exponential enters, stated without softening]\label{rem:honest}
The scoping note \cite{scoping} states the constants change by ``$M_2\mapsto
M_2e^{\Lambda\Tep}$, $M_3\mapsto M_3e^{\Lambda\Tep}$''. \textbf{The second half
of that is wrong, and this note corrects it.} Grönwall pays for the
\emph{length of the interval on which it is run}, and there are three such
intervals, of three different lengths:
\begin{center}
\begin{tabular}{lll}
\toprule
interval & what it controls & factor \\
\midrule
$[\thit,\trel]$, length $\Tep+O(\beta^{-1})$ & $M_2$, clause (ivA), (i), (ii) & $e^{\Lambda_{\rm ep}\Tep}$ \\
$[0,\thit]$, length $\thit$ & $\rho(\thit)$, subdominant & $e^{\Lambda_{\rm pre}\thit}$ \\
$[\trel,H]$, length $H-\trel$ & $M_3$, clause (ivB) & $e^{\Lambda_{\rm post}(H-\trel)}$ \\
\bottomrule
\end{tabular}
\end{center}
Only the first is an episode factor, and the first two \emph{compose}: the
pre-hit factor multiplies the initial datum \emph{inside} the bracket that the
episode factor then scales (\eqref{eq:Brate}); they are not two alternative
readings of one factor. Clause~(ivB) \emph{does} need the factor, and it needs a
\emph{horizon} one; a worst-case reading of all three is $e^{\Xmax\Cres H}$.
\textbf{And the threshold inherits the worst of them for every clause.} It is
tempting to say ``clauses (i)--(iii) only pay $e^{\Lambda\Tep}$'', and that is
true of the constants \emph{displayed in their statements}; but $\beta_0$ is one
number, the maximum over the five rows of \cite[Tab.~5]{paper}, and one of those
rows is fed by clause~(ivB). So the \emph{statements} of (i)--(iii) are
episode-clean while the \emph{threshold above which they hold} is not. We
separate the two rather than round in our favour. Two honest mitigations, neither of which we claim as a
theorem. First, the sharp form of \eqref{eq:gronwall} carries
$\exp\bigl(\int\Lambda(u)\,du\bigr)$, and by the fine reading of
\cref{rem:readings} $\Lambda\le\|\nabla U\|/\varphi_v$ decays along a converging
flow, so $\int_{\trel}^{H}\Lambda\,du$ may be far below $\Lambda_{\rm post}H$.
Second, clause~(ivB) is a statement about the \emph{tail}, and neither the
first-order coefficient \eqref{eq:release} nor clauses (i)--(iii) depend on it.
The law is not at risk; the tail constant is.
\end{remark}

%======================================================================
\section{\texorpdfstring{The coefficient: $c\mapsto c\,v_0$}{The coefficient: c -> c v0}}\label{sec:coef}

\begin{lemma}[normal form with trainable dead readout, L1$''$]\label{lem:L1}
Under \hyplabel{H4$'$} and \hyplabel{H6$'$}, the margin of $A$ on the face $j$
obeys
\begin{equation}\label{eq:margin}
\dot z=-\frac1n\sum_i \res_i\,v_0\,\sigma(\beta z_{A,i})(x_ix_j+1)
=c\,v_0(t)\,\sigma(\beta z)\,R(t)+O(e^{-\beta\delta_1}),
\end{equation}
so every occurrence of $c$ in L1$'$--L6$'$ of \cite{paper} is replaced by
$c\,v_0(t)$. Moreover $v_0$ is quasi-constant on the episode:
\begin{equation}\label{eq:v0const}
\dot v_0=-\frac1n\sum_i\res_i\,g_\beta(z_{0,i}),
\qquad
\bigl|v_0(t)-v_0(\thit)\bigr|=O(\beta^{-2}\Tep)\quad\text{on }[\thit,\trel].
\end{equation}
\end{lemma}

\begin{proof}[Proof sketch]
\eqref{eq:margin} is the chain rule with the readout carried along; the
$O(e^{-\beta\delta_1})$ collects the data $i\ne j$, whose $A$-margins are
$\le-\delta_1$ by \hyplabel{H4$'$}. For \eqref{eq:v0const}: the readout gradient
is proportional to the \emph{output} of its own unit,
$\partial L/\partial v_0=\frac1n\sum_i\res_ig_\beta(z_{0,i})$. On the face,
clause~(i) parks the margin at $z_{\rm park}=-\beta^{-1}\log\beta(1+o(1))$, and
$g_\beta(-\beta^{-1}\log\beta)=\beta^{-1}\log(1+\beta^{-1})=O(\beta^{-2})$; off
the face, $g_\beta(z_{0,i})\le\beta^{-1}e^{-\beta\delta_1}$. With
$|\res_i|\le\Cres$ this gives $|\dot v_0|=O(\beta^{-2})$ on the episode, and
integrating over $[\thit,\trel]$, of length $\Tep+O(\beta^{-1})$, gives
\eqref{eq:v0const}. \emph{This uses clause~(i), which the bootstrap L4$''$
establishes, so the argument is part of the same induction and not prior to it}
--- the same circularity, closed the same way, as in \cite[L4$'$]{paper}.
\end{proof}

\begin{remark}[the mechanism, in one sentence]
A dead unit freezes its own readout. The gradient of $v_0$ is proportional to
the output of unit $0$, which is beyond all orders precisely while the unit is
dead --- the very same fact that produces the plateau, read in the readout
coordinate instead of the margin coordinate.
\end{remark}

\begin{remark}[L5$''$, the cancellation, is untouched]\label{rem:L5}
The cancellation of \cite[L5$'$, Rem.~B.4]{paper} that produces the first-order
coefficient is arithmetic in $R$ and in the constant prefactor of the
quadrature. Replacing $c$ by $c\,v_0$ with $v_0$ constant on the episode changes
that prefactor and nothing else: the accumulated area $\Ar^{\hyb}$ and hence
$\texit$ are \textbf{unchanged}, because $v_0$ factors out of the integrand and
the exit time is defined by the \emph{zero} of the area. The arrival speed
becomes $c|v_0||R^{\hyb}(\thit)|$, the exit speed $c\,v_0R^{\hyb}(\texit)$, and
the coefficient inherits $1/v_0$. This is why the law survives.
\end{remark}

%======================================================================
\section{The Degrau 2 theorem}\label{sec:thm}

\begin{theorem}[entry--exit law, forcing form, trainable readout]\label{thm:d2}
Let $f_\theta$ be as in \eqref{eq:net} with \emph{all} parameters trained,
$A$ the distinguished unit with margin $z$ on the face $j$, and $R$ the forcing
\cite[(2)]{paper} generated by the live block $B$. Assume
\hyplabel{H0$'$}--\hyplabel{H5$'$} of \cite[\S3.1]{paper}, read on the tube, and
\hyplabel{H6$'$}. Put $\Lambda_{\rm ep},\Lambda_{\rm pre},\Lambda_{\rm post}
\le\Xmax\Cres$ as in \cref{lem:L2}. Then there is $\beta_0<\infty$ such that for
$\beta\ge\beta_0$ clauses \textup{(i)}, \textup{(ii)}, \textup{(iii)},
\textup{(ivA)}, \textup{(ivB)} of \cite[Thm.~6]{paper} hold, with exactly two
changes.

\smallskip\noindent
\textbf{(a) The first-order coefficient carries the dead readout.}
\begin{equation}\label{eq:release}
t^\beta_{\rm release}
=\texit-\frac{1}{c\,v_0(\thit)\,R^{\hyb}(\texit)}\cdot\frac1\beta+o(\beta^{-1}),
\end{equation}
$\texit$ being, as before, the first zero of the hybrid area $\Ar^{\hyb}$, which
is itself unchanged (\cref{rem:L5}). Clause~(i) reads
$u=c\,v_0(\thit)\beta|\Ar^{\hyb}(t)|(1+o(1))$.

\smallskip\noindent
\textbf{(b) The constants gain exponentials, on three intervals of three
different lengths, and a power of the guard.}
\[
M_2\ \mapsto\ M_2^{\,\prime\prime}=e^{\Lambda_{\rm ep}\Tep}\,\frac{C_fC_4}{c\,|v_0(\thit)|},
\qquad
M_3\ \mapsto\ e^{\Lambda_{\rm post}(H-\trel)}\,M_3 ,
\]
the pre-hit datum $\rho(\thit)$ acquiring $e^{\Lambda_{\rm pre}\thit}$
\emph{inside} the bracket that $e^{\Lambda_{\rm ep}\Tep}$ then multiplies
(\cref{eq:Brate}: the two pre-hit factors \emph{compose}, they are not
alternatives). The dependence on the guard is
\[
M_2^{\,\prime\prime}\ \lesssim\ \varphi_v^{-2}\,e^{\Lambda_{\rm ep}\Tep}\,M_2 ,
\]
\textbf{a second power, not a first}: one factor $|v_0|^{-1}$ is displayed, and
a second is hidden inside $C_4$, whose first term is
$2/\bigl(c|v_0|\min(|R^{\hyb}(\thit)|,R^{\hyb}(\texit))\bigr)$ (\cref{sec:unif},
U2 and U4). Finally $\beta_0$ \textbf{inherits the largest of the three
exponentials}, namely the horizon one $e^{\Xmax\Cres H}$, together with a power
$\varphi_v^{-p}$, $p\le3$ --- and it does so \emph{for every clause}, including
(i)--(iii), because $\beta_0$ is a single threshold defined as the max over the
five rows of \cite[Tab.~5]{paper} and one of those rows is fed by
clause~(ivB).

\smallskip
In particular \textbf{the exponent $\Lambda$ admits a bound free of $m$}, so
$\beta_0\lesssim O(m)\cdot\varphi_v^{-p}\,e^{\Xmax\Cres H}$ at fixed forcing and
\cite[Prop.~7]{paper} (polynomial threshold in $m$) survives at fixed
$\Xmax$, $\Cres$, $\Tep$, $H$, $\varphi_v$.
\end{theorem}

\begin{proof}
The six lemmas of \cite[App.~B]{paper} are rewritten as L1$''$--L6$''$; only
L2$'$ changes substantively, and only L1$'$ changes cosmetically.
\begin{center}
\begin{tabular}{@{}p{0.17\textwidth}p{0.77\textwidth}@{}}
\toprule
lemma & what changes \\
\midrule
L1$''$ & the factor $v_0(t)$ in the margin equation \eqref{eq:margin}; $c\mapsto cv_0$
throughout; the additive quadrature identity and its three-window bound are
algebra in $(\Psi,u,\sigma)$ and are untouched. \Cref{lem:L1}. \\
L2$''$ & \textbf{substantive.} Step~1: identity with $\bar H$ in place of $G_B$
\eqref{eq:Deltaode}. Steps~2a/2b: \emph{intact} --- they are about $A$'s output and
the live softplus defects, not about the class of $B$. Step~3: Gr\"onwall in
place of the free PSD estimate, exponent $\Lambda_{\rm ep}$ \eqref{eq:gronwall}.
Steps~4,~5: same, on their own intervals. \Cref{lem:L2}. \\
L3$''$ & entry corner: intact in form; the arrival speed is $c|v_0|\,|R^{\hyb}(\thit)|$,
positive by \hyplabel{H1$'$} and \hyplabel{H6$'$}. \\
L4$''$ & sticking/release for prescribed forcing: intact in form, with $c\mapsto cv_0$;
the closed induction now propagates $2M_2^{\,\prime\prime}$ and additionally
carries \eqref{eq:v0const}, whose proof uses clause~(i) of the same induction
(\cref{lem:L1}). \\
L5$''$ & \textbf{intact.} The cancellation is arithmetic in $R$ and in the
prefactor $cv_0$ (\cref{rem:L5}); it never sees the class of $B$. This is the
structural reason \eqref{eq:release} holds at all. \\
L6$''$ & exit corner, branch selection, scalar reduction: intact; the exit
speed is $v'=cv_0R^{\hyb}(\texit)>0$ by \hyplabel{H2$'$} and \hyplabel{H6$'$}. \\
\bottomrule
\end{tabular}
\end{center}
Given L1$''$--L6$''$, the assembly of \cite[App.~B.7]{paper} is verbatim: the
scalar reduction of L6$''$ turns the episode into the one-dimensional
entry--exit problem, L1$''$ puts it in the normal form
$d\Psi(u)/dt=-c\,v_0\beta R+c\,v_0\log(1+w)+\mathcal E$, L3$''$/L6$''$ supply
the two corners, L4$''$ the sticking law and the release instant for prescribed
$R$, and L2$''$ closes the loop by bounding the perturbation of $R$ caused by
the $\beta$-flow's own deviation. Substituting $c\mapsto cv_0$ with $v_0$
constant to $O(\beta^{-2}\Tep)$ propagates through the assembly linearly, giving
\eqref{eq:release}; the constants collect the three Gr\"onwall factors of
\cref{lem:L2}.
\end{proof}

\begin{remark}[what the theorem does \emph{not} say]
It does not remove \hyplabel{H4$'$}: the activation pattern of the live units is
still frozen. Degrau~2 is orthogonal to (O2) of \cite{paper}, not a step towards
it. It also does not claim sharpness of $\Lambda$, does not quantify the
$o(1)$ of the entry corner ((O3) of \cite{paper} is inherited unchanged), and
does not establish \eqref{eq:v0const} outside the episode.
\end{remark}

%======================================================================
\section{\texorpdfstring{The discriminating experiment: is the $v_0$ factor real?}{The discriminating experiment: is the v0 factor real?}}\label{sec:num}

\subsection{Why the previous campaign could not decide}

The scoping campaign \cite{scoping} ran the whole pipeline with trainable
readouts and $v_0(0)=1$, and reported
$\max_{43\text{ instances}}|v_0(\thit)-1|=8.7\cdot10^{-6}$. That measurement
\emph{confirms the mechanism} of \cref{lem:L1} --- the dead unit does freeze its
readout --- but it makes the two candidate predictors
\[
\text{naive}=-\frac{1}{c\,R^{\hyb}(\texit)}
\qquad\text{versus}\qquad
\text{corrected}=-\frac{1}{c\,v_0(\thit)\,R^{\hyb}(\texit)}
\]
numerically identical. It was registered, correctly, as a \emph{null of
instrument}: the experiment cannot distinguish the paper's law from
\eqref{eq:release} when $v_0\equiv1$.

\subsection{Design of the discriminating run}

Falsifiable prediction, fixed before the run: \emph{if $v_0$ freezes at its
initial value, then initializing $v_0(0)\in\{0.5,2.0\}$ makes the two predictors
differ by a factor $1/v_0\in\{2,\tfrac12\}$, and the measured
$\beta(t_{\rm release}-\texit)$ must follow the corrected one.} The condition
that kills \eqref{eq:release}: \emph{if \texttt{err\_corr} is not systematically
below \texttt{err\_naive} across the sample, item~(a) is wrong and the
coefficient of \cite[Thm.~6]{paper} stands unchanged.}

Everything else is held at the scoping values: same engine
(\path{verifica_forcing.py}), same drawing family, $m=4$ live units plus one
dead, $\beta=10^4$, $dt_0=10^{-5}$ hybrid / $2\cdot10^{-6}$ coupled, $H=1.8$,
$n=4$, same acceptance filter (\hyplabel{H1$'$}--\hyplabel{H4$'$} plus the sign
guard of \hyplabel{H6$'$} local to the window). Script:
\path{05-experimentos/t2-n-neuronios/degrau2_v0.py}; output
\path{degrau2_v0.jsonl}; $2$ workers, \texttt{nice -19}.
Only $v_0(0)$ is changed. Note that this changes the instances themselves ---
different $f$, different residuals, different episodes --- so the comparison
\emph{across} $v_0$ levels is not paired; the comparison that decides is
\texttt{err\_corr} against \texttt{err\_naive} \emph{within} each instance,
which is exactly paired.

\subsection{Result}\label{sec:result}

% degrau2-numeros.tex -- gerado a mao a partir de degrau2_v0.log / degrau2_v0.jsonl
% (05-experimentos/t2-n-neuronios), run de 2026-07-28, 10.3 min, 2 workers nice -19.
% [RESULT 2026-07-28]

\paragraph{The readout does freeze, at both levels.}
Over the accepted instances,
\[
\max\bigl|v_0(\thit)-v_0(0)\bigr|
=1.2\cdot10^{-6}\ \ (v_0(0)=0.5),
\qquad
3.6\cdot10^{-6}\ \ (v_0(0)=2.0),
\]
so \eqref{eq:v0const} holds far more tightly than the $O(\beta^{-2}\Tep)$ it
promises, and the two predictors now differ by a clean factor $2$ and
$\tfrac12$ respectively. The instrument is no longer null.

\paragraph{The corrected law wins, and by the predicted factor.}
Acceptance: $22/30$ at $v_0(0)=0.5$ and $26/30$ at $v_0(0)=2.0$ (rejections: $11$
$\beta$-runs with no release inside the horizon, $1$ death before $t_{\min}$; no
instance was rejected by the sign guard of \hyplabel{H6$'$}). Writing
$\mu=\beta(t^\beta_{\rm release}-\texit)$ for the measured coefficient:

\begin{center}
\begin{tabular}{@{}lcccc@{}}
\toprule
& \multicolumn{2}{c}{$v_0(0)=0.5$ ($N=22$)} & \multicolumn{2}{c}{$v_0(0)=2.0$ ($N=26$)} \\
\cmidrule(lr){2-3}\cmidrule(lr){4-5}
& $\mu/\text{naive}$ & $\mu/\text{corr.}$ & $\mu/\text{naive}$ & $\mu/\text{corr.}$ \\
\midrule
predicted by \eqref{eq:release} & $2.000$ & $1.000$ & $0.500$ & $1.000$ \\
\textbf{measured, median}       & $\mathbf{2.049}$ & $\mathbf{1.024}$ & $\mathbf{0.496}$ & $\mathbf{0.991}$ \\
measured, range                 & $[1.63,2.95]$ & $[0.81,1.47]$ & $[0.48,0.69]$ & $[0.97,1.38]$ \\
\bottomrule
\end{tabular}
\end{center}

\noindent
The medians of the ratio to the \emph{naive} predictor are $2.049$ and $0.496$
against the arithmetic values $2$ and $\tfrac12$ that \eqref{eq:release}
predicts, with no free parameter and no fit. Relative errors:

\begin{center}
\begin{tabular}{@{}lccc@{}}
\toprule
 & median & p90 & max \\
\midrule
$|\mu-\text{naive}|/|\text{naive}|$, pooled $N=48$ & $51.4\%$ & --- & --- \\
$|\mu-\text{corrected}|/|\text{corrected}|$, pooled $N=48$ & $\mathbf{2.76\%}$ & $17.3\%$ & $47.3\%$ \\
\ \ at $v_0(0)=0.5$ & $5.80\%$ & $16.7\%$ & $47.3\%$ \\
\ \ at $v_0(0)=2.0$ & $2.19\%$ & $15.3\%$ & $37.7\%$ \\
\bottomrule
\end{tabular}
\end{center}

\noindent
\textbf{Paired verdict: the corrected predictor beats the naive one in $47$ of
$48$ instances}, with median ratio $\texttt{err\_corr}/\texttt{err\_naive}=0.043$
--- a factor $23$. The single exception is at $v_0(0)=2.0$ and is an instance
whose corrected error is already the sample maximum ($37.7\%$), i.e.\ an
instance where the residual $\beta_0$ effect dominates the $v_0$ effect, not one
where the naive law is right.

\smallskip\noindent
\emph{The falsification condition fixed before the run was not met.} Item~(a) of
\cref{thm:d2} stands: the first-order coefficient carries the dead unit's
readout.

\paragraph{The lemma is again saturated, and Gr\"onwall is again not the driver.}
Measuring $L_B=\max(0,-\lambda_{\min}(\nabla^2U))$ by exact eigenvalues of the
$12\times12$ Hessian along the hybrid trajectory, against
$\Lambda=\max_k\|c_k\|$: $\max_{48}|L_B-\Lambda|=7.1\cdot10^{-15}$ (relative
$3.0\cdot10^{-15}$). Ranges: $\Lambda\in[2.36,7.07]$ (median $4.34$),
$\Tep\in[0.0145,0.351]$ (median $0.0626$), and the episode Gr\"onwall factor
\[
e^{\Lambda\Tep}\in[1.08,\ 2.87],\qquad\text{median }1.28 .
\]
The correlation of $\log\texttt{err\_corr}$ with $\log e^{\Lambda\Tep}$ is
$\mathbf{-0.45}$ --- \emph{negative}, as in \cite{scoping}, so the exponential
factor cannot be the causal channel for the residual error --- while the
correlation with $\log(1/R^{\hyb}(\texit))$ is $+0.29$, the same
loss-of-transversality driver that (O4) of \cite{paper} identifies for the
frozen readout. Both correlations are weaker than in \cite{scoping}
($-0.63$, $+0.60$), which is expected: pooling two $v_0$ levels mixes two
different drawing families. \emph{We do not claim a decomposition;} $N=48$ at a
single $\beta$ excludes Gr\"onwall as the dominant channel by the sign, and
nothing more.

\paragraph{Guard and balancedness.}
The empirical guard is $\varphi_v=\min_k\min_t|v_k|=5.0\cdot10^{-3}$ over the
whole sample (median over instances $0.106$); no sign crossing occurred inside a
window in any accepted instance. Balancedness drift
$\max_k|B_k(t)-B_k(0)|$ over the episode: $\le6.0\cdot10^{-5}$ on the hybrid
($dt=10^{-5}$) and $\le9.5\cdot10^{-6}$ on the $\beta$-run
($dt=2\cdot10^{-6}$), a ratio $\approx6.3$ against a step ratio of $5$ --- the
drift is Euler error, not softplus defect, as \cref{sec:open}\,(P3) predicts.


%======================================================================
\section{\texorpdfstring{Uniformity in $m$, and the constants}{Uniformity in m, and the constants}}\label{sec:unif}

We track what changes relative to \cite[Rem.~B.9]{paper}.

\begin{itemize}
\item[(U1)] \textbf{The first-order coefficient is still exactly $m$-free.} It
is $-1/(c\,v_0(\thit)R^{\hyb}(\texit))$: built from $c$ (a property of the datum
$j$ and of $n$), from \emph{one} readout, and from the hybrid forcing at one
instant. $m$ cannot appear.
\item[(U2)] $C_4$ gains $1/|v_0|$ through the arrival and exit speeds,
\[C_4=\frac{2}{c\,|v_0|\min\bigl(|R^{\hyb}(\thit)|,R^{\hyb}(\texit)\bigr)}+\frac{\Tep}{a_0},\]
and therefore $C_4\le$ its frozen-readout value times $\varphi_v^{-1}$ by
\hyplabel{H6$'$}. Still a functional of the forcing alone, given $v_0$.
\item[(U3)] $C_f=\sup_t\|\nabla_{\theta_B}f_j\|/n$ now has three components per
unit instead of two, namely
$\bigl(v_k\mathbf1_{k,j}x_j,\ v_k\mathbf1_{k,j},\ \mathbf1_{k,j}z_{k,j}\bigr)$,
so that
\[
C_f\ \le\ \frac1n\Bigl[\Xmax\bigl(\textstyle\sum_kv_k^2\bigr)^{1/2}
+\bigl(\textstyle\sum_kz_{k,j}^2\bigr)^{1/2}\Bigr],
\]
still $O(\sqrt m)$ for both readout scalings. \emph{Unchanged in order.}
\item[(U4)] $M_2^{\,\prime\prime}=e^{\Lambda_{\rm ep}\Tep}C_fC_4/(c|v_0|)
\lesssim\sqrt m\,\varphi_v^{-2}e^{\Xmax\Cres\Tep}$. Compare \cite{paper}:
$M_2\lesssim\sqrt m$; and \cite[L2]{paper}, the two-neuron form:
$\sqrt m\,e^{\Theta(mH)}$. \textbf{The trainable readout costs one exponential
whose exponent admits a bound free of the width (the realized exponent is not yet measured, (P7)); the $m\cdot H$ restriction does not come
back.} This is the whole content of \cref{lem:curv} at the level of constants.
\item[(U5)] The live-unit term \eqref{eq:Brate}, third group, is unchanged:
$C_5\sqrt m\,\Cres e^{-\beta\delta_B}$, affecting $\beta_1$ logarithmically and
$\beta_2\propto m$ linearly. The $3m$ versus $2m$ dimension count changes
nothing in order.
\item[(U6)] \textbf{New: $\varphi_v$.} The guard enters $\beta_0$ in three
places, all through speeds and all as $\varphi_v^{-1}$: the arrival speed $v$ of
L3$''$, the exit speed $v'$ of L6$''$ and $C_A$, and $C_4$ via (U2). It does
\emph{not} enter $\Lambda$ (\cref{rem:readings}, crude reading). So
$\beta_0\lesssim O(m)\,\varphi_v^{-p}e^{\Xmax\Cres H}$ for a small integer $p$
($p\le3$ by inspection of the five rows of \cite[Tab.~5]{paper}); we do not
multiply the rows out, for the reason given in (O1) of \cite{paper}.
\item[(U7)] \textbf{The horizon, and what $\beta_0$ inherits.} The worst of the
three exponentials is $e^{\Lambda_{\rm post}(H-\trel)}$ in $M_3$: it is
clause~(ivB) that carries a horizon factor, and the constants displayed in
clauses (i)--(iii) and (ivA) carry only $e^{\Lambda_{\rm ep}\Tep}$, with the
pre-hit factor composed inside. \textbf{But $\beta_0$ is a single threshold and
inherits $e^{\Xmax\Cres H}$ for all clauses alike} (\cref{rem:honest}): the
statements are episode-clean, the threshold is not.
\end{itemize}

%======================================================================
\section{Open steps, stated in the Degrau 1 register}\label{sec:open}

\begin{itemize}
\item[(P1)] \textbf{The comparison with the frozen-readout baseline is not
paired.} The scoping campaign reports median error $1.96\%$ ($n=43$) against
$0.80\%$ for the frozen readout ($n=23$), a factor $\approx2.5$. The two rows
are different instances: with trainable readouts this drawing family yields
episodes $3$--$10\times$ shorter ($\Tep$ median $0.058$ against $0.2$--$0.8$).
The factor is indicative, not a contrast. A paired design would freeze the
instance and toggle only the trainability of $(v_k)_{k\ge1}$; it has not been
run.
\item[(P2)] \textbf{$N$ is small, and $\beta$ is a single value.} Every number
in \cref{sec:num} and in \cite{scoping} comes from $m=4$, $n=4$, one $\beta$
($10^4$), one integrator (explicit Euler), and samples of a few tens of
instances ($N=48$ here, $N=43$ there). There is no $\beta$-sweep, no $m>4$, and no discretization control
for the $v_0$ campaign. In particular, the claim that the residual error is the
$\beta_0(\text{forcing})$ of Degrau~1 rather than a failure of the law rests on
the sign of a correlation, not on a collapse under a $\beta$-sweep. The
$\beta$-sweep on the worst instances is the cheapest next test and it is not
done.
\item[(P3)] \textbf{The balancedness route to \hyplabel{H6$'$} is open, but it
is now a \emph{one-unit} problem.} Under ReLU gradient flow
$B_k=v_k^2-\|(w_k,b_k)\|^2$ is exactly conserved \cite{duhulee}; under softplus
the defect is beyond all orders inside \hyplabel{H4$'$} (measured
$2.8\cdot10^{-5}$, $1.7\cdot10^{-7}$, $\sim10^{-13}$ at
$\beta=10,10^2,10^3$). With \hyplabel{H6$'$} reduced to the dead unit
(\cref{rem:H6only0}) the sufficient condition is $B_0(0)>0$, on \emph{one}
triple, giving $\varphi_v=\sqrt{B_0(0)}$.

This is a genuine simplification of the obstruction. In the $m+1$-unit version
the balanced rescaling $(v,w,b)\mapsto(sv,w/s,b/s)$ --- which does not change
$f$ --- removes $100\%$ of observed sign crossings ($0/150$ against $19/300$)
but breaks \hyplabel{H4$'$} by shrinking the live margins, and there is a
genuine competition. Applied to unit $0$ \emph{alone}, the rescaling touches no
live margin at all; the only constraint it must respect is that $A$'s own
margins still satisfy the $\delta_1$ clause of \hyplabel{H4$'$} and that $A$
still dies in the window. Whether a rescaling factor $s$ exists that meets both
has \textbf{not} been checked, numerically or otherwise. The campaigns here use
the initialization of \cite{paper}, which has $B_0(0)<0$, so we \emph{measure}
$\varphi_v$ rather than construct it.
\item[(P4)] \textbf{Saturation of \cref{lem:curv} is unexplained.} See
\cref{rem:sat}. The inequality is proved; the observed equality is not.
\item[(P10)] \textbf{$v_0=0$ is a genuine singularity; $v_k=0$ is not.} The
coefficient \eqref{eq:release} blows up as $v_0(\thit)\to0$, and this is not an
artefact of the estimate: at $v_0=0$ the dead unit is functionally absent, the
margin equation \eqref{eq:margin} has $\dot z\equiv0$ to leading order, and
there is no release to predict. The degradation is graceful and of exactly the
same species as (O1) of \cite{paper} --- a forcing-side quantity in a
denominator, with no lower bound available from any initialization criterion.
A \emph{live} readout crossing zero, by contrast, is innocuous
(\cref{rem:H6only0}): no step consumes a lower bound on $v_k$, $k\ge1$. The
acceptance filter of the campaigns rejects both, which is stricter than
necessary and costs sample size, not validity.
\item[(P5)] \textbf{The averaged-Hessian bound is taken on the tube.} $\Lambda_{\rm ep}$
in \eqref{eq:Lamsup} is a supremum over the segment between the two
trajectories, and is therefore bounded through \hyplabel{H5$'$}-on-the-tube.
Since $\|\Delta\|=O(\log\beta/\beta^2)$, the tube is thin and the distinction is
immaterial in practice; but the bootstrap that makes it rigorous is the same
continuity argument as in \cite[L4$'$]{paper} and we have not rewritten it.
\item[(P7)] \textbf{The $m$-sweep of $\Lambda$ is a pilot, not a measurement.}
The claim that survives proof is that the \emph{bound} $\Xmax\Cres$ is $m$-free;
the claim that the \emph{realized} $\Lambda$ is flat in $m$ rests on $N=4,4,1$
instances at $m=4,8,16$ (\cref{sec:lema}). The design that would settle it, and
which we have not run: $N\approx20$ accepted instances per level,
$m\in\{4,8,16,32\}$, same drawing family and the fine $dt$, reporting the median
and the interquartile range of $\Lambda$ and of $\lambda_{\max}(G_B)$ per level,
with the pre-registered prediction that the first is flat and the second linear.
Note that $m=32$ needs a $96\times96$ Hessian per sample and a coarser
\texttt{HESS\_EVERY}; cost is a few CPU-hours at two workers.
\item[(P8)] \textbf{We measure $\Lambda$, not the quantity the estimate
actually needs.} Step~3 of \cref{lem:L2} uses the Rayleigh quotient of the
averaged Hessian \emph{along the realized $\Delta$},
\[
q(t)\ :=\ -\frac{\Delta^\top\bar H(t)\Delta}{\|\Delta\|^2}\ \le\ \Lambda_{\rm ep},
\]
and the Gr\"onwall factor that the argument truly pays is
$\exp\bigl(\int q\,dt\bigr)$, not $e^{\Lambda\Tep}$. The two coincide only if
$\Delta$ aligns with the minimizing eigenvector. \emph{Measuring $q(t)$ is the
experiment that separates the two live explanations of \cref{sec:num}}: if
$q\approx\Lambda$, the exponent is right and the estimate is tight, and the
negative correlation of the residual error with $e^{\Lambda\Tep}$ is telling us
the error lives elsewhere; if $q\ll\Lambda$, the Gr\"onwall step is simply loose
and $\Lambda$ is an artefact of the worst-case direction. The instrumentation is
cheap --- $\Delta$ and $\bar H$ are both already computed in
\path{degrau2_readout.py} --- and it has not been done.
\item[(P9)] \textbf{Selection is conditional on the treatment.} See
\cref{sec:num}: the acceptance filter rejected $8/30$ at $v_0(0)=0.5$ against
$4/30$ at $v_0(0)=2.0$, so the two arms are not samples of the same population
and the between-arm comparison of error levels is confounded with selection.
The within-instance paired comparison, which is what decides item~(a), is
untouched by this. Quantifying the selection --- e.g.\ by re-running with the
horizon $H$ enlarged until the rejection rates match --- is not done.
\item[(P6)] \textbf{Inherited unchanged from \cite{paper}}: (O2) frozen live
pattern, (O3) the unquantified $o(1)$ of the entry corner, (O1) the forcing
dependence of $\beta_0$ (every row degenerates as $a_0\to0$ or
$R^{\hyb}(\texit)\to0$, and the numerics of \cref{sec:num} see exactly that).
\end{itemize}

\bigskip
\noindent\textbf{Epistemic labels.}
\Cref{lem:curv} and \cref{lem:L1}\,\eqref{eq:margin}: \emph{[ESTABLISHED]}
(closed computation, verified symbolically and numerically).
Identity \eqref{eq:Deltaode}: \emph{[ESTABLISHED]} (mean value theorem for
gradients).
Balancedness: \emph{[ESTABLISHED]}, Du--Hu--Lee \cite{duhulee}, source opened
and read 2026-07-27.
\Cref{eq:v0const}: \emph{[SPECULATION with derivation]} --- the computation is
done, the rigorous proof needs the L4$''$ bootstrap.
\Cref{thm:d2} as a whole: \emph{[SPECULATION]} --- proof sketch at lemma
granularity, no step audited by a referee.
Numbers of \cref{sec:num}: \emph{[RESULT 2026-07-27]}.
Saturation of \cref{lem:curv} on trajectories: \emph{[SPECULATION]} for the
mechanism, \emph{[RESULT]} for the measurement.

\begin{thebibliography}{9}
\bibitem{paper} L.~Caloi Santos, \emph{An entry--exit law for a dying softplus
unit} (\texttt{06-notas/paper/paper-v4.tex}), 2026. Theorem~6 (forcing form),
\S3.1; App.~B (proof); App.~B.9 (uniformity); App.~B.10 (open items (O1)--(O6)).
\bibitem{scoping} \emph{Degrau 2 --- o teorema de forcing com readout
trein\'avel (scoping)}, \texttt{06-notas/2026-07-27-degrau2-scoping.md}, 2026-07-27.
\bibitem{duhulee} S.~S.~Du, W.~Hu, J.~D.~Lee, \emph{Algorithmic regularization
in learning deep homogeneous models: layers are automatically balanced},
NeurIPS 2018, arXiv:1806.00900.
\end{thebibliography}

\end{document}
